% Source: https://artofproblemsolving.com/community/c3h1405032_ant_on_icosahedron
% Stable identity: aops-topic-1405032-q01
Call the vertices $T$ (top) and $B$ (bottom). The icosahedron consists of the $5$ faces around $T$, a band of $10$ faces between the upper ring of $5$ vertices $U_i$ and the lower ring of $5$ vertices $L_i$, and the $5$ faces around $B$. Unfold the faces $TU_1U_2$, $U_1U_2L_1$, $U_2L_1L_2$, $L_1L_2B$ into the plane; they form part of a triangular lattice of unit triangles:
\[
T=(0,0),\ U_{1,2}=\left(\mp\dfrac{1}{2},-\dfrac{\sqrt3}{2}\right),\ L_1=(0,-\sqrt3),\ L_2=(1,-\sqrt3),\ B=\left(\dfrac{1}{2},-\dfrac{3\sqrt3}{2}\right).
\]

\begin{asy}
usepackage("amsmath");
size(4.5cm);
defaultpen(fontsize(10pt));
real h=sqrt(3)/2;
pair T=(0,0),U1=(-0.5,-h),U2=(0.5,-h),L1=(0,-2*h),L2=(1,-2*h),Bv=(0.5,-3*h);
pair Ua=(-1.5,-h),Ub=(1.5,-h),La=(-1,-2*h),Lb=(2,-2*h);
fill(T--U1--U2--cycle,gray(0.9)); fill(U1--U2--L1--cycle,gray(0.9)); fill(U2--L1--L2--cycle,gray(0.9)); fill(L1--L2--Bv--cycle,gray(0.9));
draw(T--U1--U2--cycle); draw(U1--L1--U2--L2--L1); draw(L1--Bv--L2);
draw(T--Bv,red+linewidth(1));
dot(T);dot(Bv);
label("$T$",T,N); label("$B$",Bv,S); label("$U_1$",U1,W); label("$U_2$",U2,E); label("$L_1$",L1,W); label("$L_2$",L2,E);
label("$\sqrt7$",(0.51,-1.52*h),red+fontsize(9pt));
label("\mbox{}",(min(currentpicture,user=true)+max(currentpicture,user=true))/2,black); // keeps the text after the figure black when LaTeX typesets the labels inline
\end{asy}

The straight segment $TB$ in this net stays inside the four shaded faces (it crosses $U_1U_2$ at $x=\dfrac{1}{6}$, then $U_2L_1$, then $L_1L_2$ at $x=\dfrac{1}{3}$), so it is a path on the surface of length
\[
TB=\sqrt{\left(\dfrac{1}{2}\right)^2+\left(\dfrac{3\sqrt3}{2}\right)^2}=\sqrt{\dfrac{1}{4}+\dfrac{27}{4}}=\sqrt7 .
\]
\textbf{No path is shorter.} A shortest path that avoids vertices is a straight segment in the unfolding of the faces it crosses, and in such an unfolding all vertices lie on the triangular lattice, with lattice neighbours being icosahedron neighbours. The lattice points at distance less than $\sqrt7$ from $T$ are at distance $1$, $\sqrt3$ or $2$, and these are reached from $T$ in at most $2$ lattice steps, so they are images of vertices at most $2$ edges from $T$; but $B$ is $3$ edges from $T$. So such a segment has length at least $\sqrt7$. A path through a vertex $V\ne T,B$ is at least (distance $T$ to $V$) $+$ (distance $V$ to $B$), and one of these vertices is at least $2$ edges from $T$ or from $B$, which gives length at least $1+\sqrt3>\sqrt7$.

Hence the answer is $\boxed{\sqrt7}$. (The same unfolding idea gives $\sqrt3$ for the octahedron and $\sqrt5$ for the cube.)
